\subsection{Evolución del Plan - ¿Cómo cambia el plan cuando cambia el proyecto?}\label{sec:spmp-6-6}
El PMP relaciona alcance, datos, simulaciones FDS y recursos con las actividades de administración del proyecto. La adaptación de los estándares ISO/IEC/IEEE 29148~\cite{std29148} e ISO/IEC/IEEE 12207~\cite{std12207} contempla conservar cambios, decisiones y versiones junto con sus efectos y evidencias. El historial de la sección \ref{sec:spmp-1} identifica las modificaciones documentales realizadas.

La gestión de una modificación comenzará con la identificación de la necesidad de cambio y la descripción del ajuste propuesto; a continuación, se evaluará su efecto sobre el alcance, los entregables, los requisitos, los recursos, la calendarización y el presupuesto del proyecto. Este análisis se relacionará con la administración de requerimientos descrita en la sección \ref{sec:spmp-10-1}.\textcolor{black}{La propuesta se discutirá entre los integrantes del equipo y el director del proyecto. La autoridad para aprobar los cambios corresponde al director y a todos los integrantes en conjunto. La aprobación precederá a la incorporación de la modificación al plan. }\textcolor{black}{Una vez aprobado el cambio, el integrante Daniel Ramírez será responsable de incorporarlo en los apartados afectados del PMP} y se revisará la coherencia de sus referencias internas y de los documentos relacionados. \textcolor{black}{Los cambios se realizarán mediante el servidor documental y las herramientas descritas en la sección \ref{sec:spmp-8-2}. Para cada actualización se realizará un commit que indique la versión del documento a la que se actualiza y describa los cambios respecto a la versión anterior. Además, los cambios se registrarán en el historial de cambios del documento, en la sección \ref{sec:spmp-1}. La versión entregada al final del periodo 2026-3 se establecerá como la versión 1.0 del PMP.}

El cambio aprobado se comunicará al equipo y a las personas que participen en las actividades afectadas, y su implementación se reflejará en el trabajo y en los documentos correspondientes, y se vinculará con el seguimiento del progreso de la sección \ref{sec:spmp-10-2} para comprobar que lo realizado corresponda con la decisión aprobada.

Los acuerdos de convivencia y toma de decisiones del equipo se describen en la sección \ref{sec:spmp-12-1}.
